\newpage\handout[true]
{\textsc{Math 2106-B: Foundations of Mathematical Proof}}{Fall 2025}
{Homework 3}
{\textit{Prof.} \href{https://ceheitsch.github.io/webpage/}{Dr. Christine Heitsch}}{Due Date: September 18th}
{Math 2106-B: HW3}
{\large
  \noindent
  \textbf{Name:} \HandoutAuthor \smallskip \\

  \noindent
  \textbf{GTID:} 903743506 \smallskip \\

  \noindent
  %%% List anyone with whom you discussed the problems
  \textbf{Collaborators:} None. This material reflects independent preparation. \smallskip \\

  \noindent
  %%% List all resources used OTHER than the textbook, lecture notes, 
  %%% webwork problems, ICA questions, and previous homework assignments.
  \textbf{Outside resources:} This work was informed by MATH 2106 course materials 
  (ICAs: in-class activities and WebWork contents) from \HandoutInstitution\ (the course textbook is 
  Chartrand \parencite{chartrand_mathematical_2018}).   
  Outside references specifically include
  textbooks written by 
  Grimaldi \parencite{grimaldi_discrete_2004} (used in MATH 3012) and 
  Rosen \parencite{rosen_discrete_2019} (used in CS 2050). \smallskip \\

  \noindent
  \textbf{Notice}: The answers contain cross-references throughout this document as clickable hyperlinks in most PDF viewers. However, the Gradescope PDF view might not support clickable hyperlinks.
}\newpage
% \printbibliography[
  heading=bibintoc,   
  title={References}     % change “References” text if needed
]
% `heading=bibnumbered` if you prefer it numbered chapter (in \tableofcontents)
% `heading=bibintoc`    if you prefer it unnumbered but still in the ToC\newpage
% The following table presents the standard logical equivalences from these sources together with 
the author's own (or the student, myself, Jaehoon Song) clarifications and labels where applicable.


\begin{table}[h!]
  \centering
  \renewcommand{\arraystretch}{1.2}
  \setlength{\tabcolsep}{8pt}
  \begin{tabular}{|p{5.0cm}|p{5.0cm}|p{5.2cm}|}
    \hline
    \rowcolor[HTML]{E0E0E0}
    \multicolumn{2}{|l|}{\textbf{Logical Equivalences}} & \textbf{Name} \\
    \hline
    $\,p \AND q \equiv q \AND p$ & $\,p \OR q \equiv q \OR p$ & Commutative Law \\
    \hline
    $\, (p \AND q) \AND r \equiv p \AND (q \AND r)$ & $\, (p \OR q) \OR r \equiv p \OR (q \OR r)$ & Associative Law \\
    \hline
    $\,p \AND (q \OR r) \equiv (p \AND q) \OR (p \AND r)$ & $\,p \OR (q \AND r) \equiv (p \OR q) \AND (p \OR r)$ & Distributive Law \\
    \hline
    $\,\overline{p \AND q} \equiv \sNOT{p} \OR \sNOT{q}$ & $\,\overline{p \OR q} \equiv \sNOT{p} \AND \sNOT{q}$ & DeMorgan's Law \\
    \hline
    \multicolumn{2}{|l|}{$\,\overline{\overline{p}} \equiv p$} & Double Negation law \\
    \hline
    $\,p \AND \boldsymbol{T} \equiv p$ & $\,p \OR \boldsymbol{F} \equiv p$ & Identity Law \\
    \hline
    $\,p \OR \boldsymbol{T} \equiv \boldsymbol{T}$ & $\,p \AND \boldsymbol{F} \equiv \boldsymbol{F}$ & Domination Law \\
    \hline
    $\,p \AND p \equiv p$ & $\,p \OR p \equiv p$ & Idempotent Law \\
    \hline
    \multicolumn{2}{|l|}{$\,p \THEN q \equiv \sNOT{p} \OR q$} & Implication Equivalence \\
    \hline
    \multicolumn{2}{|l|}{$\,\overline{p \THEN q} \equiv p \AND \sNOT{q}$} & Negation of Implication \\
    \hline
    \multicolumn{2}{|l|}{$\,p \THEN q \equiv \sNOT{q} \THEN \sNOT{p}$} & Contrapositive Equivalence \\
    \hline
    \multicolumn{2}{|l|}{$\,p \OR \sNOT{p} \equiv \boldsymbol{T}$} & Law of Tautology \\
    \hline
    \multicolumn{2}{|l|}{$\,p \AND \sNOT{p} \equiv \boldsymbol{F}$} & Law of Contradiction \\
    \hline
  \end{tabular}
\end{table}\newpage
% 
The following definitions are presented from standard mathematical sources (\textbf{Sets}).
Clarifications and labels have been added where applicable.

\begin{define}{Basic Set Concepts}\phantomsection\label{def:basic-set-concepts} % \hyperref[def:basic-set-concepts]{YOUR_CONTEXT}
  \begin{enumerate}
    \item \textbf{Universal Set:} A universal set $\setUNIV$ is a
    set that contains all objects under consideration in a particular context or problem.
    \item \textbf{Empty Set:} The empty set, denoted $\setNULL$,
    is the unique set that contains no elements.
    \item \textbf{Subset:} Let $A$ and $B$ be sets. $A$ is a
    subset of $B$ ($A \subseteq B$) if (and only if)
    for all $x \in A$, $x \in B$.
    \item \textbf{Set Equality:} Let $A$ and $B$ be sets. The sets $A$ and $B$ are
    equal ($A = B$) if $A \subseteq B$ and $B \subseteq A$.
    \item \textbf{Power Set:} Let $A$ be a set. The power set of $A$, denoted $\sPOWER{A}$ or $2^A$,
    is the set of all subsets of $A$. That is, $\mathcal{P}(A) = \{S\in\setUNIV : S \subseteq A\}$.
  \end{enumerate}
\end{define}
The following definitions extend the understanding of set operations
and properties.


\begin{define}{Set Operations}\phantomsection\label{def:set-operations} % \hyperref[def:set-operations]{YOUR_CONTEXT}
  \begin{enumerate}
    \item \textbf{Complementary Set:} Let $A$ be a subset of the universal set $\setUNIV$.
    The complement of $A$, denoted $A^c$ or $\overline{A}$, is the set of all elements
    in $\setUNIV$ that are not in $A$. That is, $A^c = \{x \in \setUNIV : x \notin A\}$.
    \item \textbf{Union:} Let $A$ and $B$ be sets. The union of $A$ and $B$,
    denoted $A \cup B$, is the set of all elements that are in $A$, or in $B$, or in both.
    That is, $A \cup B = \{x : x \in A \lor x \in B\}$.
    \item \textbf{Intersection:} Let $A$ and $B$ be sets. The intersection of $A$ and $B$,
    denoted $A \cap B$, is the set of all elements that are in both $A$ and $B$.
    That is, $A \cap B = \{x : x \in A \land x \in B\}$.
    \item \textbf{Set Difference:} Let $A$, $B$ be subsets of a universal set $\setUNIV$.
    The set difference $A \sMINUS B$ is defined to be the set
    $\{ x \in \setUNIV \LDIV x \in A \text{ and } x \notin B\}$. \vspace*{10px}\\
    \textit{It is sometimes called the relative complement of $B$ in $A$,
    and denoted in the textbook as $A - B$.}
    \item \textbf{Cartesian Product:} Let $A$ and $B$ be sets. The Cartesian product of
    $A$ and $B$, denoted $A \sTIMES B$, is the set consisting of all ordered pairs:
    \[
    A \sTIMES B = \{(a, b) \LDIV a \in A, b \in B\}.
    \]
  \end{enumerate}
\end{define}


\newpage

The following claims are given in the reference as well. Let $A,B\in U$, a universal set.


\begin{claim}{Subset of Union: $A \subseteq \sOR{A}{B}$}\phantomsection\label{claim:subset-union} % \hyperref[claim:subset-union]{YOUR_CONTEXT}
  \begin{subprove}[bluecolor]
      Let $x\in U$, assume $x\in A$. Then the statement \inlinequote{$x\in A$ or $x\in B$} is true.
      Hence, by the definition, $\sOR{A}{B}=\{x\in U\LDIV x\in A$ or $x\in B\}$, $x\in \sOR{A}{B}$.
      Thus, $A \subseteq \sOR{A}{B}$.
  \end{subprove}
\end{claim}

\begin{claim}{Empty Subset: $\setNULL \subseteq A$}\phantomsection\label{claim:empty-subset} % \hyperref[claim:empty-subset]{YOUR_CONTEXT}
  \begin{subprove}[bluecolor]
      Let $x\in U$, assume $x\in \setNULL$. But there are no elements in $\setNULL$, so
      $x\in \setNULL$ is false. Since the premise is false, the implication if $x\in \setNULL$,
      then $x\in A$ is true. Thus, $\setNULL \subseteq A$ by definition of \hyperref[def:basic-set-concepts]{subsets}.
  \end{subprove}
\end{claim}

\begin{claim}{Subset as Equality: $\sOR{A}{B}=A$ if and only if $B\subseteq A$}\phantomsection\label{claim:subset-equality} % \hyperref[claim:subset-equality]{YOUR_CONTEXT}
  \begin{subprove}[bluecolor]
      Let $x\in U$, first, we assume $B\subseteq A$, we have already shown
      $A \subseteq \sOR{A}{B}$ by the claim, \hyperref[claim:subset-union]{Subset of Union}. 
      To show $\sOR{A}{B} \subseteq A$, assume $x\in \sOR{A}{B}$.
      By definition, $x\in A$ or $x\in B$. Considering $x\in B$ and assumption
      $B \subseteq A$, $x\in B$ implies $x\in A$. Thus, $\sOR{A}{B} \subseteq A$.
      Hence, if $B\subseteq A$, then $\sOR{A}{B} = A$. \vspace*{10px}\\
      Now, we assume $\sOR{A}{B}=A$. To show $B\subseteq A$, we now assume $x\in B$. Since $x\in B$, by definition
      of union, $x\in \sOR{A}{B} = A$. Therefore, $B\subseteq A$.
  \end{subprove}
\end{claim}


%%%%%%%%%%%%%%%%%%% 9/8/25

Let $\setUNIV$ be a set. 
\begin{claim}{ICA6}\phantomsection\label{claim:ica6} % \hyperref[claim:ica6]{YOUR_CONTEXT}
  For all $A,B\in\setUNIV$, $$A\sTIMES B =\setNULL\THEN \group{A=\setNULL\OR B=\setNULL}$$
  \begin{subprove}[bluecolor]
    Let $A,B\in\setUNIV$. We prove $A\sTIMES B=\setNULL$. Assume it is not the case that 
    $A=\setNULL$ or $B=\setNULL$. Thus, we have $A=\setNULL$ and $B=\setNULL$. 
    Since $A\ne\setNULL$ by assumption, there exists an element $a\in A$. Likewise, there 
    is some $b\in B$ since $B\ne\setNULL$. Thus, $\group{a,b}\in\sBUILD{(x,y)}{x\in A,y\in B}$ 
    by definition. But then, $A\sTIMES B\ne\setNULL$ by definition. Hence, the claim holds. 
  \end{subprove}
\end{claim}


\begin{define}{Relatively Prime Integers (Coprime)}\phantomsection\label{def:relatively-prime} % \hyperref[def:relatively-prime]{YOUR_CONTEXT}
  Let $x,y \in \mathbb{Z}$. The integers $x$ and $y$ are said to be
  \textbf{relatively prime} if and only if
  $\gcd(x,y) = 1$.
\end{define}
\newpage
% The following definitions and axioms are presented from standard mathematical sources (\textbf{Number Theory}) together with the author's own (or the student, myself, Jaehoon Song) clarifications and labels where applicable.

\begin{define}{Integers}[def:integers][red]
  $\setZ$: the set of all integers
\end{define}

\begin{define}{Rational Numbers}[def:rational-numbers][red]
  $\setQ$: the set of all rational numbers
\end{define}

\begin{define}{Real Numbers}[def:real-numbers][red]
  $\setR$: the set of all real numbers
\end{define}

\begin{define}{Natural Numbers}[def:natural-numbers][red]
  $\setN$ or $\setZ_{>0}$: the set of all natural numbers which are the set of positive integers (not including $0$)
\end{define}

\begin{define}{Positive Real Numbers}[def:positive-real-numbers][red]
  $\setR_{>0}$: the set of all positive real numbers
\end{define}

\begin{define}{Non-negative Real Numbers}[def:non-negative-real-numbers][red]
  $\setR_{\geq 0}$: the set of all non-negative real numbers
\end{define}


\newpage

The following definitions and claims extend the understanding of integers.

Here are the ICA2 Definitions and Axioms as follows.

\begin{define}{Even Integer}[def:even-integer][red]
  An integer $n$ is \underbar{even} if there exists 
  (existential quantification logic $\exists$) an integer $k$ such that
  \[
  n = 2k.
  \]
\end{define}

\begin{define}{Odd Integer}[def:odd-integer][red]
  An integer $n$ is \underbar{odd} if 
  there exists some integer $k$ such that
  \[
  n = 2k + 1.
  \]
\end{define}

\begin{define}{Closure under Addition}[def:closure-addition][red]
  The sum of two integers is again an 
  integer. (Closure under addition)
\end{define}

\begin{define}{Closure under Multiplication}[def:closure-multiplication][red]
  The product of two integers is again 
  an integer. (Closure under multiplication)
\end{define}

\begin{define}{Odd if and only if Not Even}[def:odd-iff-not-even][red]
  An integer is odd if and only if it 
  is not even.
\end{define}

Here are the ICA3 Definitions and Axioms as follows.

\begin{define}{Universal Closure under Addition}[def:universal-closure-addition][red]
  The sum of two integers (Universal quantification for all, for every, for each) is again an integer. (Closure under addition)
\end{define}

\begin{define}{Universal Closure under Multiplication}[def:universal-closure-multiplication][red]
  The product of two integers (Universal quantification for all, for every, for each) is again an integer. (Closure under multiplication)
\end{define}

\begin{define}{Even Integer Square}[def:even-integer-square][red]
  An integer $n$ is even if and only if $n^2$ is even.
\end{define}

\newpage

The following claims are given in the reference as well.

\begin{define}{Sum of Even Integers}[claim:sum-even-integers][red]
  The sum of two even integers is again an even integer.
  \begin{subprove}[red]
    Let $x$ be an integer and $y$ be an integer ($\forall x,y \in \setZ$). 
    Suppose $x$ is even and $y$ is even. Then, by \nameref{def:even-integer}, there exists an integer 
    $k$ and an integer $j$ such that $x=2k$ and $y=2j$. Then, 
    $x+y=\left(2k\right)+\left(2j\right)=2\left(k+j\right)$. By \nameref{def:closure-addition}, since $k$ and 
    $j$ are integers, then so is $\left(k+j\right)$. Hence, by definition, $x+y$ is an 
    even integer.
  \end{subprove}
\end{define}

The following definitions extend the understanding of integers including divisibility.

\begin{define}{Divisibility}[def:divisibility][red]
  An integer $n$ divides an integer $m$, denoted $n \LDIV m$, if there 
  exists an integer $k$ such that $m=n k$. Otherwise, $n \nmid m$.
\end{define}

The following definitions extend the classification of real numbers into rational and irrational numbers.

\begin{define}{Rational Number}[def:rational-number][red]
  A real number $x$ is \underbar{rational} if there exist integers $a$ and $b$ with $b \neq 0$ such that $bx = a$.
  Otherwise, $x$ is \underbar{irrational}.

  \textbf{Note}: Compare the rational number definition with the even/odd axiom \inlinequote{odd iff not even} from \nameref{def:odd-iff-not-even}.
\end{define}

The following claims are given in the reference as well.

\begin{define}{Irrationality of $\sqrt{2}$}[claim:sqrt2-irrational][red]
  $\sqrt{2}$ is irrational.
  \begin{subprove}[red]
    Suppose $\sqrt{2}$ is rational. Then, by \nameref{def:rational-number}, there exists $a,b\in \setZ$ 
    with $b\ne 0$ such that $b\cdot\sqrt{2}=a$. We may further assume that any factors 
    common to $a$ and $b$ have already been removed.\\

    Then, squaring both sides, this yields $2\cdot b^2 = a^2$.
    Since products of integers are again integers by \nameref{def:closure-multiplication}, we have that $a^2,b^2\in \setZ$.
    Thus, by definition, $a^2$ is even. Therefore, by \nameref{def:even-integer-square}, $a$ is even. 
    Thus, there exists $c\in\setZ$ such that $a=2c$. Then
    $$
    2b^2=a^2={2c}^2=4c^2
    $$
    Then, $b^2=2c^2$. Then, as before, $b^2$ is even, so $b$ is even.
    However, $a$ and $b$ both being even contradicts the fact that any factors 
    common to $a$ and $b$ have already been removed.
  \end{subprove}
\end{define}






%%%%%%%%%%%%%%%%%%% New style (refactor model)
\begin{define}{Relatively Prime Integers (Coprime)}[def:relatively-prime][red]
  Let $x,y \in \mathbb{Z}$. The integers $x$ and $y$ are said to be
  \textbf{relatively prime} if and only if
  $\gcd(x,y) = 1$.
\end{define}\newpage






%%%%%%%%%%%%%%%%%%%%%%%%%%%%%%%%%%%%%%%%%%%%%%%%%%%%%%%%%%%%%%%%%%%%%%%
% References
%%%%%%%%%%%%%%%%%%%%%%%%%%%%%%%%%%%%%%%%%%%%%%%%%%%%%%%%%%%%%%%%%%%%%%%
\textbf{Remark}: The following definitions (or results) are already proved in class or on previous ICA, 
Wbwk, Hmwk assignments, and in Sections 4.4, 4.5, and 4.6 of the textbook.
\begin{define}{Set Equality}[def:set-equality][red]
  Let $A$ and $B$ be sets. The sets $A$ and $B$ are equal ($A = B$) iff $A \subseteq B$ and $B \subseteq A$.
\end{define}
\begin{define}{Subset Relation}[def:subset][red]
  Let $A$ and $B$ be sets. $A$ is a subset of $B$ ($A \subseteq B$) iff for all $x \in A$ only if $x \in B$.
\end{define}  
\begin{define}{Union of Sets}[def:union][red]
  Let $A$ and $B$ be sets. The union of $A$ and $B$, denoted $A \cup B$, is the 
  set of all elements that are in $A$, or in $B$, or in both. 
  That is, $A \cup B = \{x : x \in A \text{ or } x \in B\}$.
\end{define}
\begin{define}{Intersection of Sets}[def:intersection][red]
  Let $A$ and $B$ be sets. The intersection of $A$ and $B$, denoted $A \cap B$, is the 
  set of all elements that are in both $A$ and $B$. That is, $A \cap B = \{x : x \in A \text{ and } x \in B\}$.
\end{define}
\begin{define}{Set Difference}[def:set-difference][red]
  Let $A$, $B$ be subsets of a universal set $\setUNIV$. The set difference $A \sMINUS B$ is defined to be the set $\{ x \in \setUNIV \LDIV x \in A \text{ and } x \notin B\}$. \vspace*{5px}\\
  \textit{It is sometimes called the relative complement of $B$ in $A$, and denoted in the textbook as $A - B$.}
\end{define}
\begin{define}{Power Set}[def:power-set][red]
  Let $A$ be a set. The power set of $A$, denoted $\sPOWER{A}$ or $2^{\abs{A}}$, is the set of all subsets of $A$. That is, $\sPOWER{A} = \{S\in\setUNIV : S \subseteq A\}$.
\end{define}
\begin{define}{Complementary Set}[def:complementary-set][red]
  Let $A$ be a subset of the universal set $\setUNIV$. The complement of $A$, 
  denoted $A^c$ or $\sNOT{A}$, is the set of all elements in $\setUNIV$ that are not in $A$. 
  That is, $\sNOT{A} = \{x \in \setUNIV : x \notin A\}$.
\end{define}
\begin{define}{Cartesian Product}[def:cartesian-product][red]
  Let $A$ and $B$ be sets. The Cartesian product of $A$ and $B$, denoted $A \sTIMES B$, is the set consisting of all ordered pairs:
  \[
  A \sTIMES B = \{(a, b) \LDIV a \in A, b \in B\}.
  \]
\end{define}
\begin{define}{Pairwise Disjoint Sets}[def:pairwise-disjoint][red]
  Let $A_1, A_2, \ldots, A_n$ be sets. These sets are \textbf{pairwise disjoint} if for all $i \neq j$, we have $A_i \cap A_j = \emptyset$.
\end{define}


% Functions
\newpage

\begin{define}{Function}[def:function][red]
  Let $A$ and $B$ be nonempty sets. A \textbf{function} $f$ from $A$ to $B$, denoted $f : A \to B$, 
  is a relation, $f\in A\sTIMES B$ such that 
  \begin{enumerate}
    \item For every $a \in A$, there exists a $b \in B$ such that $f(a) = b$.
    \item For every $a \in A$ and for all $b, c \in B$, if $f(a) = b$ and $f(a) = c$, then $b = c$.
    \item The set $A$ is called the \textbf{domain} of $f$, 
    and the set $B$ is called the \textbf{codomain} of $f$.
    \item If $f(a)=b$, then $b$ is called the \textbf{image} of $a$ under $f$.
    \item The function maps $a$ to $b$, and $f$ is often called a mapping.
  \end{enumerate}
\end{define}

\begin{define}{Domain and Codomain}[def:domain-codomain][red]
  If $f : A \to B$ is a function, then $A$ is called the \textbf{domain} of $f$, and $B$ is called the \textbf{codomain} of $f$.
\end{define}

\begin{define}{Image and Inverse Image}[def:image][red]
  Let $f : A \to B$ be a function, $C\subseteq A$, and $D\subseteq B$. 
  The \textbf{image} $f(C)$ is the set $\{f(c) \in B : c \in C\} \subseteq B$.
  The \textbf{inverse image} $f^{-1}(D)$ is the set $\{a \in A : f(a) \in D\} \subseteq A$.
\end{define}

\begin{define}{Injective, Surjective, Bijective Functions}[def:injectivity-surjectivity-bijective][red]
  Let $f : A \to B$ be a function.
  \begin{itemize}
    \item $f$ is \textbf{injective} (one-to-one) iff for all $a_1, a_2 \in A$, if $f(a_1) = f(a_2)$, then $a_1 = a_2$.
    \item $f$ is \textbf{surjective} (onto) iff for every $b \in B$, there exists $a \in A$ such that $f(a) = b$.
    \item $f$ is \textbf{bijective} iff it is both injective and surjective.
  \end{itemize}
\end{define}


\begin{define}{Theorem}[def:bijective][red]
  The inverse relation $f^{-1}$ is a function if and only if $f$ is both injective and surjective.
  % \begin{subprove}
  %   ($\Rightarrow$) Assume $f^{-1}$ is a function. To show $f$ is injective, let $a_1, a_2 \in A$ such that $f(a_1) = f(a_2)$. 
  %   Then, $(a_1, f(a_1)) \in f$ and $(a_2, f(a_2)) \in f$. Since $f^{-1}$ is a function, $(f(a_1), a_1) \in f^{-1}$ and $(f(a_2), a_2) \in f^{-1}$. 
  %   By the definition of a function, $a_1 = a_2$. Thus, $f$ is injective.
    
  %   To show $f$ is surjective, let $b \in B$. Since $f^{-1}$ is a function, there exists an $a \in A$ such that $(b, a) \in f^{-1}$. 
  %   By the definition of inverse relation, $(a, b) \in f$, so $f(a) = b$. Thus, $f$ is surjective.
    
  %   ($\Leftarrow$) Assume $f$ is both injective and surjective. To show $f^{-1}$ is a function, let $(b, a_1), (b, a_2) \in f^{-1}$. 
  %   Then, $(a_1, b), (a_2, b) \in f$. Since $f$ is injective, it follows that $a_1 = a_2$. Thus, $f^{-1}$ satisfies the definition of a function.
  % \end{subprove}

  \begin{subprove}
    We consider the inverse relation $f^{-1} \subseteq B \times A$ for the function 
    \[
    f : A \to B \quad \text{on nonempty sets $A,B$.}
    \]
    Assume $f$ is a bijection, i.e., both one-to-one and onto. Let $b \in B$. Since $f$ is onto, by definition, there exists $a \in A$ with $f(a) = b$. So $(a,b) \in f$ and $(b,a) \in f^{-1}$. Hence for all $b \in B$ there exists $a \in A$ such that $(b,a) \in f^{-1}$. 

    Now let $a, a' \in A$. Suppose $(b,a) \in f^{-1}$ and $(b,a') \in f^{-1}$. Then $(a,b), (a',b) \in f$, so $f(a) = b$ and $f(a') = b$. But $f$ is one-to-one by assumption, and so $a = a'$. Hence we have shown that each $b \in B$ has at least one image under the inverse relation $f^{-1}$ and no more than one. Thus the inverse relation is a function whenever $f$ is a bijection.

    \medskip
    Conversely, \dots
  \end{subprove}
\end{define}






%%%%%%%%%%%%%%%%%%%%%%%%%%%%%%%%%%%%%%%%%%%%%%%%%%%%%%%%%%%%%%%%%%%%%%%
% Problems
%%%%%%%%%%%%%%%%%%%%%%%%%%%%%%%%%%%%%%%%%%%%%%%%%%%%%%%%%%%%%%%%%%%%%%%
\begin{enumerate} 
  \newpage
  \item Let $A$, $B$, $C$, and $D$ be subsets of a universal set $\setUNIV$. 
  Prove the following propositions.~\footnote{
    \textbf{Notice}: The answers contain cross-references throughout this document as clickable hyperlinks in most PDF viewers. 
    However, the Gradescope PDF view might not support clickable hyperlinks.
  }\\
  \textit{This proves that the three subsets are a partition of $(A \sOR B)$. Partitions will be very important to understanding equivalence relations.}
  \begin{enumerate}
    \item $(A \sOR B)=(A \sAND B) \sOR(A \sMINUS B) \sOR(B \sMINUS A)$.~\footnote{
      I intentionally avoided repetitive or excessive cross-references for better readability.
    }
    \begin{subprove}
      By definition of \nameref{def:set-equality}, $(A \sOR B)=(A \sAND B) \sOR(A \sMINUS B) \sOR(B \sMINUS A)$ 
      requires showing both directions of \nameref{def:subset}s, as follows.\vspace*{5px}\\
      ($\subseteq$) Let $x \in A \sOR B$ be arbitrary. Then, $x \in A$ or $x \in B$ 
      by definition of \nameref{def:union}. Consider the following cases.
      \begin{itemize}
        \item If $x \in A$ and $x \in B$, then $x \in A \sAND B$ by definition of \nameref{def:intersection}.
        \item If $x \in A$ and $x \notin B$, then $x \in A \sMINUS B$ by definition of \nameref{def:set-difference}.
        \item If $x \notin A$ and $x \in B$, then $x \in B \sMINUS A$ by definition.
      \end{itemize}
      In all cases, $x \in (A \sAND B) \sOR(A \sMINUS B) \sOR(B \sMINUS A)$. 
      Hence, proof by cases, $(A \sOR B) \subseteq (A \sAND B) \sOR(A \sMINUS B) \sOR(B \sMINUS A)$.\vspace*{5px}\\
      ($\supseteq$) Let $x \in (A \sAND B) \sOR(A \sMINUS B) \sOR(B \sMINUS A)$ be arbitrary. 
      Then, $x \in A \sAND B$ or $x \in A \sMINUS B$ or $x \in B \sMINUS A$. Consider the following cases.
      \begin{itemize}
        \item If $x \in A \sAND B$, then $x \in A$ and $x \in B$ by definition, so $x \in A \sOR B$.
        \item If $x \in A \sMINUS B$, then $x \in A$ by definition, so $x \in A \sOR B$ by definition.
        \item If $x \in B \sMINUS A$, then $x \in B$ by definition, so $x \in A \sOR B$ by definition.
      \end{itemize}
      Hence, proof by cases, $(A \sAND B) \sOR(A \sMINUS B) \sOR(B \sMINUS A) \subseteq A \sOR B$.\vspace*{5px}\\
      Therefore, by ($\subseteq$) and ($\supseteq$), 
      $(A \sOR B)=(A \sAND B) \sOR(A \sMINUS B) \sOR(B \sMINUS A)$.
    \end{subprove}

    \item $(A \sAND B)$, $(A \sMINUS B)$, and $(B \sMINUS A)$ are all pairwise disjoint.
    \begin{subprove}
      It is required to show that any two of these sets have empty intersection,
      \begin{itemize}
        \item To claim $(A \sAND B) \sAND (A \sMINUS B) = \setNULL$,
        assume there exists $x \in (A \sAND B) \sAND (A \sMINUS B)$ for 
        the sake of contradiction. Then, $x \in A \sAND B$ and $x \in A \sMINUS B$. 
        Since $x \in A \sAND B$, it follows that $x \in A$ and $x \in B$ by definition of \nameref{def:intersection}. 
        Since $x \in A \sMINUS B$, it follows that $x \in A$ and $x \notin B$ by definition of \nameref{def:set-difference}, 
        which contradicts $x \in B$. Therefore, $(A \sAND B) \sAND (A \sMINUS B) = \setNULL$.
        \item To claim $(A \sAND B) \sAND (B \sMINUS A) = \setNULL$, without loss of generality, 
        it is symmetrical to the previous case, with the roles of sets $A$ and $B$ swapped.
        Hence, trivially, $(A \sAND B) \sAND (B \sMINUS A) = \setNULL$.
        \item To claim $(A \sMINUS B) \sAND (B \sMINUS A) = \setNULL$,
        assume there exists $x \in (A \sMINUS B) \sAND (B \sMINUS A)$ for 
        the sake of contradiction. Then, $x \in A \sMINUS B$ and $x \in B \sMINUS A$. 
        Since $x \in A \sMINUS B$, it follows that $x \in A$ and $x \notin B$ by definition of \nameref{def:set-difference}. 
        Since $x \in B \sMINUS A$, it follows that $x \in B$ and $x \notin A$ by definition of \nameref{def:set-difference}, 
        which contradicts both $x \notin B$ and $x \in A$ at the same time. 
        Therefore, $(A \sMINUS B) \sAND (B \sMINUS A) = \setNULL$.
      \end{itemize}
      Therefore, $(A \sAND B)$, $(A \sMINUS B)$, and $(B \sMINUS A)$ are all pairwise disjoint by definition of \nameref{def:pairwise-disjoint}.
    \end{subprove}
  \end{enumerate}
  




  \newpage
  \item Let $A$, $B$, $C$, and $D$ be subsets of a universal set $\setUNIV$. 
  Prove or disprove the following propositions.~\footnote{
    \textbf{Notice}: The answers contain cross-references throughout this document as clickable hyperlinks in most PDF viewers. 
    However, the Gradescope PDF view might not support clickable hyperlinks.
  }\\
  \textit{Likewise, Cartesian product of sets (above) and the power set of a set (below) will be very important in the next part of the course.}
  \begin{enumerate}
    \item $(A \sTIMES B) \sAND(C \sTIMES D) \subseteq(A \sAND C) \sTIMES(B \sAND D)$.
    \begin{subprove}
      Let $(x,y) \in (A \sTIMES B) \sAND(C \sTIMES D)$ be arbitrary for the sake of directly showing ($\subseteq$). 
      Then, $(x,y) \in A \sTIMES B$ and $(x,y) \in C \sTIMES D$. 
      Since $(x,y) \in A \sTIMES B$, it follows that $x \in A$ 
      and $y \in B$ by definition of \nameref{def:cartesian-product}. Since $(x,y) \in C \sTIMES D$, it follows 
      that $x \in C$ and $y \in D$ as well.
      Thus, $x \in A \sAND C$ 
      and $y \in B \sAND D$ by definition of \nameref{def:intersection}, 
      so $(x,y) \in (A \sAND C) \sTIMES(B \sAND D)$.
      Hence, $(A \sTIMES B) \sAND(C \sTIMES D) \subseteq(A \sAND C) \sTIMES(B \sAND D)$.
    \end{subprove}

    \item $(A \sAND C) \sTIMES(B \sAND D) \subseteq(A \sTIMES B) \sAND(C \sTIMES D)$.
    \begin{subprove}
      Let $(x,y) \in (A \sAND C) \sTIMES(B \sAND D)$ be arbitrary for the sake of directly showing ($\subseteq$). 
      Then, $x \in A \sAND C$ and $y \in B \sAND D$ by definition of \nameref{def:cartesian-product}. 
      Since $x \in A \sAND C$, it follows that $x \in A$ and $x \in C$ by definition of \nameref{def:intersection}. 
      Since $y \in B \sAND D$, it follows that $y \in B$ and $y \in D$.
      Therefore, $(x,y) \in A \sTIMES B$ and $(x,y) \in C \sTIMES D$ since logical AND (\AND) is commutative.
      Thus, $(x,y) \in (A \sTIMES B) \sAND(C \sTIMES D)$.
      Hence, $(A \sAND C) \sTIMES(B \sAND D) \subseteq(A \sTIMES B) \sAND(C \sTIMES D)$.
    \end{subprove}

    \item $(A \sTIMES B) \sOR(C \sTIMES D) \subseteq(A \sOR C) \sTIMES(B \sOR D)$.
    \begin{subprove}
      Let $(x,y) \in (A \sTIMES B) \sOR(C \sTIMES D)$ be arbitrary for the sake of directly showing ($\subseteq$). 
      Then, $(x,y) \in A \sTIMES B$ or $(x,y) \in C \sTIMES D$ by definition of \nameref{def:union}. 
      \begin{itemize}
        \item If $(x,y) \in A \sTIMES B$, then $x \in A$ and $y \in B$ by definition of \nameref{def:cartesian-product}. \label{prob:proof-2}
        Since $x \in A$, it follows that $x \in A \sOR C$ by definition of \nameref{def:union}. 
        Since $y \in B$, it follows that $y \in B \sOR D$ by definition of \nameref{def:union}. 
        Therefore, $(x,y) \in (A \sOR C) \sTIMES(B \sOR D)$.
        \item If $(x,y) \in C \sTIMES D$, without loss of generality, it is symmetrical to the previous case.
        Therefore, it trivially implies $(x,y) \in (A \sOR C) \sTIMES(B \sOR D)$.
      \end{itemize}
      In both cases, $(x,y) \in (A \sOR C) \sTIMES(B \sOR D)$. 
      Hence, proof by cases, $(A \sTIMES B) \sOR(C \sTIMES D) \subseteq(A \sOR C) \sTIMES(B \sOR D)$.
    \end{subprove}

    \item $(A \sOR C) \sTIMES(B \sOR D) \subseteq(A \sTIMES B) \sOR(C \sTIMES D)$.
    \begin{subdisprove}
      Let $A=\{1\}, B=\{2\}, C=\{3\}, D=\{4\}$ be subsets of a universal set $\setUNIV$ 
      for the sake of counterexample. 
      Then, $(A \sOR C) \sTIMES(B \sOR D)=\{1,3\} \sTIMES \{2,4\}=\{(1,2),(1,4),(3,2),(3,4)\}$ 
      and $(A \sTIMES B) \sOR(C \sTIMES D)=\{(1,2)\} \sOR \{(3,4)\}=\{(1,2),(3,4)\}$. 
      Since $(1,4) \in (A \sOR C) \sTIMES(B \sOR D)$ but $(1,4) \notin (A \sTIMES B) \sOR(C \sTIMES D)$, 
      it follows that $(A \sOR C) \sTIMES(B \sOR D) \not\subseteq (A \sTIMES B) \sOR(C \sTIMES D)$. 
      Therefore, by counterexample, the statement is \textbf{false}.
    \end{subdisprove}
  \end{enumerate}





  \newpage
  \item Let $A$, $B$ be subsets of a universal set $\setUNIV$. 
  Prove or disprove the following propositions.~\footnote{
    \textbf{Notice}: The answers contain cross-references throughout this document as clickable hyperlinks in most PDF viewers. 
    However, the Gradescope PDF view might not support clickable hyperlinks.
  }
  \begin{enumerate}
    \item $\sPOWER(A \sAND B) \subseteq \sPOWER{A} \sAND \sPOWER{B}$.
    \begin{subprove}
      Let $X \subseteq\setUNIV$. Assume $X\in\sPOWER(A \sAND B)$ be arbitrary for the sake of directly showing ($\subseteq$). 
      Then, $X \subseteq A \sAND B$ by definition of \nameref{def:power-set}. 
      Since $X \subseteq A \sAND B$, it follows that 
      $X \subseteq A$ and $X \subseteq B$ by definition of \nameref{def:intersection}. Therefore, $X \in \sPOWER{A}$ 
      and $X \in \sPOWER{B}$, so $X \in \sPOWER{A} \sAND \sPOWER{B}$.
      Hence, $\sPOWER(A \sAND B) \subseteq \sPOWER{A} \sAND \sPOWER{B}$.
    \end{subprove}
    
    \item $\sPOWER{A} \sAND \sPOWER{B} \subseteq \sPOWER(A \sAND B)$.
    \begin{subprove}
      Let $X \subseteq\setUNIV$. Assume $X\in\sPOWER{A} \sAND \sPOWER{B}$ be arbitrary for the sake of directly showing ($\subseteq$). 
      Then, $X \in \sPOWER{A}$ and $X \in \sPOWER{B}$ by definition of \nameref{def:intersection}. 
      Since $X \in \sPOWER{A}$, it follows that $X \subseteq A$ by definition of \nameref{def:power-set}. 
      Without loss of generality, (trivial) $X \subseteq B$.
      Since $X \subseteq A$ and $X \subseteq B$, it follows that $X \subseteq A \sAND B$.
      Therefore, $X \in \sPOWER(A \sAND B)$.
      Hence, $\sPOWER{A} \sAND \sPOWER{B} \subseteq \sPOWER(A \sAND B)$.
    \end{subprove}

    \item $\sPOWER(A \sOR B) \subseteq \sPOWER{A} \sOR \sPOWER{B}$.
    \begin{subdisprove}
      Let $A=\{1\}, B=\{2\}$ be subsets of a universal set $\setUNIV$ 
      for the sake of counterexample. 
      Then, $\sPOWER(A \sOR B)=\sPOWER(\{1,2\})=\{\setNULL,\{1\},\{2\},\{1,2\}\}$ by definition of \nameref{def:power-set}
      and $\sPOWER{A} \sOR \sPOWER{B}=\{\setNULL,\{1\}\} \sOR \{\setNULL,\{2\}\}=\{\setNULL,\{1\},\{2\}\}$. 
      Since $\{1,2\} \in \sPOWER(A \sOR B)$ but $\{1,2\} \notin \sPOWER{A} \sOR \sPOWER{B}$, 
      it follows that $\sPOWER(A \sOR B) \not\subseteq \sPOWER{A} \sOR \sPOWER{B}$. 
      Therefore, by the counterexample, the statement is \textbf{false}.
    \end{subdisprove}

    \item $\sPOWER{A} \sOR \sPOWER{B} \subseteq \sPOWER(A \sOR B)$.
    \begin{subprove}
      Let $X \subseteq\setUNIV$. Assume $X\in\sPOWER{A} \sOR \sPOWER{B}$ be arbitrary for the sake of directly showing ($\subseteq$). 
      Then, $X \in \sPOWER{A}$ or $X \in \sPOWER{B}$ by definition of \nameref{def:union}. 
      \begin{define}{Subset of Union: $A \subseteq \sOR{A}{B}$}[claim:subset-union][red]
        \begin{subprove}[red]
            Let $x\in U$, assume $x\in A$. Then the statement \inlinequote{$x\in A$ or $x\in B$} is true.
            Hence, by definition of \nameref{def:union}, $\sOR{A}{B}=\{x\in U\LDIV x\in A$ or $x\in B\}$, $x\in \sOR{A}{B}$.
            Thus, $A \subseteq \sOR{A}{B}$.
        \end{subprove}
      \end{define}
      If $X \in \sPOWER{A}$, then $X \subseteq A \subseteq A \sOR B$ by definition of \nameref{def:power-set} and the short claim above, \nameref{claim:subset-union}, 
      so $X \in \sPOWER(A \sOR B)$.
      Without loss of generality, the case of $X \in \sPOWER{B}$ trivially implies $X \in \sPOWER(A \sOR B)$ as well. 
      In both cases, $X \in \sPOWER(A \sOR B)$. 
      Hence, proof by cases, $\sPOWER{A} \sOR \sPOWER{B} \subseteq \sPOWER(A \sOR B)$.
    \end{subprove}
  \end{enumerate}
\end{enumerate}